\chapter{Propuesta} % aqui hay que nombrar al sistema :P
\label{cap:propuesta}

\section{Consideraciones Iniciales}

Este capítulo detalla el diseño y la implementación de nuestro 
sistema para ayudar a los analistas a explorar procesos complejos 
documentados en texto. Nuestro sistema busca resolver tres retos 
principales. Primero, la dificultad de reconstruir la narrativa 
de un caso a partir de documentos fragmentados. Segundo, la 
validación de los datos extraídos. Tercero, la integración de 
las dimensiones temporal, espacial y narrativa en una sola 
interfaz. Para ello, definimos primero los objetivos de diseño 
y las tareas analíticas que guían la propuesta. Luego, explicamos 
la implementación del sistema dividida en tres etapas clave, 
extracción, agregación y exploración interactiva.

A lo largo del capítulo usamos tres conceptos centrales. Un 
\textbf{caso} es la unidad de análisis del sistema, una 
organización, proyecto u obra sobre la que se dispone de 
múltiples documentos. Un \textbf{informe} es cada uno de esos 
documentos, que registra el estado del caso en un punto 
particular de su evolución. Una \textbf{variable} es cualquier 
atributo del caso que el sistema extrae y rastrea a lo largo 
de los informes, como el presupuesto ejecutado, los hallazgos 
identificados o el tipo de acción registrada.

\section{Objetivos y Tareas}

A partir de la revisión de la literatura y del análisis de los 
problemas identificados previamente, se definieron los siguientes 
objetivos de diseño y tareas analíticas que guían el desarrollo 
del sistema.

\subsection{Objetivos de Diseño}

\textbf{OD1. Analizar la evolución, ubicación y variables de 
los casos.} El análisis requiere examinar cómo cambian los casos 
a lo largo del tiempo, identificando sus ubicaciones geográficas 
en distintos niveles y comparando sus variables clave a partir 
de la información extraída de múltiples informes.

\textbf{OD2. Identificar similitudes y agrupaciones entre casos.}
El análisis requiere examinar el conjunto de casos para detectar 
cuáles comparten características comunes, e identificar qué tan 
atípico o recurrente es un caso dentro del conjunto.

\textbf{OD3. Verificar el origen de cada dato extraído.}
La confiabilidad del análisis requiere examinar que cada dato 
extraído proviene del documento original, permitiendo identificar 
posibles errores.

\subsection{Tareas Analíticas}

\textbf{T1. Reconstruir la evolución temporal de los casos.}
El sistema permitirá organizar cronológicamente los informes de 
cada caso para visualizar su desarrollo en el tiempo. \textbf{[OD1]}

\textbf{T2. Comparar el progreso entre casos.}
El sistema permitirá analizar la evolución de varios casos en 
paralelo, comparando sus informes y variables para identificar 
similitudes o diferencias en su desarrollo. \textbf{[OD1]}

\textbf{T3. Localizar los casos geográficamente.} 
El sistema permitirá visualizar la ubicación de los casos a 
nivel nacional, regional y local para examinar su concentración 
y detalle en diferentes niveles. \textbf{[OD1]}

\textbf{T4. Explorar casos y hallazgos similares.} 
El sistema permitirá seleccionar un caso o hallazgo para 
identificar otros elementos que compartan características o 
narrativas similares dentro del conjunto analizado. \textbf{[OD2]}

\textbf{T5. Consultar el documento original de un dato extraído.} 
El sistema permitirá acceder al fragmento del documento original 
para navegar por sus páginas y verificar el contexto exacto del 
que se extrajo cada dato. \textbf{[OD3]}

\section{Arquitectura del sistema}

\begin{figure}[htbp]
    \centering
    \includegraphics[width=\textwidth]{PIPElin.png} 
    \caption{Diseño de la arquitectura del sistema, organizado 
    en tres fases, extracción, agregación y exploración 
    interactiva.}
    \label{fig:Pipeline}
\end{figure}

Tal como se observa en la Figura~\ref{fig:Pipeline}, el sistema 
se organiza en tres fases. En la fase de \textbf{extracción}, 
cada informe se procesa de forma individual para identificar 
variables, actores, relaciones y ubicaciones, preservando la 
referencia al documento original. Los datos extraídos pasan a la 
fase de \textbf{agregación}, donde los informes de un mismo caso 
se agrupan, sus variables se sincronizan en el tiempo y se 
calculan las agregaciones geográficas. Finalmente, la fase de 
\textbf{exploración interactiva} expone esta representación al 
analista mediante vistas coordinadas.

\subsection{Extracción}

La fase de extracción procesa cada informe de forma individual 
y produce una representación estructurada de su contenido junto 
con anclajes precisos al documento original. Se compone de tres 
etapas, preprocesamiento, extracción semántica y mitigación de 
alucinaciones.

\subsubsection{Preprocesamiento}

El preprocesamiento convierte cada informe en una representación 
textual semiestructurada en la que cada elemento queda acompañado 
de su ubicación exacta en el documento, expresada como 
identificador, página y coordenadas. Optamos por una 
representación semiestructurada antes que por una estructura 
rígida porque los informes provienen de fuentes heterogéneas, y 
una plantilla fija habría requerido reglas específicas por tipo 
de documento, lo que comprometería la generalidad del 
sistema~\cite{cattan2021crossdoc}. Los informes pueden contener 
texto nativo o haber sido digitalizados mediante escaneo. En 
ambos casos utilizamos OpenDataLoader 
\ac{PDF}~\cite{opendataloader2024}, que aplica extracción directa al 
texto nativo y activa \ac{OCR} cuando lo detecta como necesario. A 
partir de esta salida se genera una representación simplificada 
que conserva solo el identificador y el texto de cada elemento, 
reduciendo el volumen de información enviada al modelo en la 
etapa siguiente.

\subsubsection{Extracción semántica}

La extracción semántica se realiza en dos consultas sucesivas 
al modelo. En la primera, se describe en lenguaje natural 
la información relevante del informe, incluyendo título, 
descripción breve, fechas, actores y ubicación geográfica. 
Para cada variable, clasificada según los seis tipos definidos 
en la Tabla~\ref{tab:tipos_variables}, el modelo devuelve tres 
elementos. Primero, el contenido de la variable. Segundo, un 
fragmento del texto que sirve de evidencia. Tercero, una 
justificación breve en lenguaje natural que explica cómo 
interpretó el fragmento para llegar a ese contenido. En la 
segunda consulta, el modelo organiza esa respuesta en la 
estructura de salida del sistema y conserva el identificador 
del elemento del que se tomó cada dato.

La decisión de exigir una justificación junto a cada dato 
responde a un argumento documentado en la 
literatura~\cite{yeh2025storyribbons}, donde se reporta que la 
capacidad de los \acs{LLM} para producir justificaciones de sus 
salidas ayuda a calibrar la confianza del usuario en la 
extracción. La justificación cumple un rol técnico adicional. 
Reduce el riesgo de respuestas no ancladas al texto, ya que 
obliga al modelo a articular el razonamiento como parte de la 
respuesta antes de comprometerse con un valor, en línea con el 
principio de razonamiento encadenado~\cite{wei2022}. Además, provee 
al analista un punto de partida para auditar el caso, sin 
tener que reconstruir desde cero por qué el modelo eligió un 
valor sobre otro.

La separación en dos consultas responde a una decisión 
adicional. Pedirle al modelo que razone sobre el contenido y 
al mismo tiempo lo serialice en un esquema rígido degrada la 
calidad de ambas tareas, según evidencia reportada en la 
literatura sobre extracción 
estructurada~\cite{li2024simple}. Al separar razonamiento y 
formato, la primera consulta se aproxima a un razonamiento 
encadenado, donde el modelo expresa libremente su análisis 
junto con la justificación. La segunda se limita a una tarea 
de transformación, donde el riesgo de alucinación es menor 
porque el contenido ya fue producido y solo resta darle 
estructura.

La tipología de seis tipos de variables que adoptamos refleja 
las distintas formas en que un dato puede evolucionar entre 
informes. Algunas variables cambian de estado, otras avanzan 
hacia una meta numérica, otras describen relaciones, y otras 
verifican una condición. Distinguirlas explícitamente permite 
que la fase de agregación trate cada tipo con la lógica 
apropiada y que la verificación elija la estrategia adecuada 
para cada dato.



\begin{table}[t]
\centering
\caption{Tipos de variables del modelo de datos.}
\label{tab:tipos_variables}
\renewcommand{\arraystretch}{1.5}
\begin{tabular}{@{}p{2.1cm} p{4.5cm} p{4cm} p{3.2cm}@{}}
\toprule
\textbf{Tipo} & \textbf{Descripción} & \textbf{Detalle interno} 
& \textbf{Ejemplo} \\
\midrule
\rowcolor{rowgray}
Seguimiento  & Datos que evolucionan entre informes, pudiendo 
resolverse, persistir o generar nuevas consecuencias & Texto, 
estado y fuente & Situaciones adversas detectadas \\
Progreso     & Valores numéricos que avanzan hacia una meta, 
comparados con un estimado por informe & Valor observado, meta, 
valor estimado y fuente & Porcentaje de presupuesto ejecutado \\
\rowcolor{rowgray}
Estado       & Categoría que describe la condición de algo en 
un informe específico & Valor categórico, nota y fuente & 
Tipo de acción (correctiva, preventiva) \\
Relacional   & Actores que participan en el caso y los roles 
que desempeñan & Lista de actores con nombre, rol y fuente & 
Organismos fiscalizados y supervisores \\
\rowcolor{rowgray}
Narrativo    & Texto estructurado con título y contenido 
extenso & Título, cuerpo del texto y fuente & Conclusiones y 
recomendaciones formales \\
Verificación & Condición booleana respaldada por evidencia y 
nivel de cumplimiento & Valor booleano, evidencia y fuente & 
Plazo de respuesta cumplido \\
\bottomrule
\end{tabular}
\end{table}

Cada dato extraído se almacena junto a tres anclajes que lo 
conectan con el documento original. La fuente, que indica de 
dónde proviene mediante el identificador del elemento, la 
página y las coordenadas. El fragmento de evidencia, que 
contiene el texto exacto desde el que se realizó la extracción. 
Y la justificación, que articula la interpretación del modelo 
sobre ese fragmento. Estos tres anclajes cumplen funciones 
complementarias. Permiten al analista consultar el documento 
desde la interfaz, habilitan la verificación automática de la 
etapa siguiente, y sirven de evidencia explícita para que el 
analista pueda calibrar su confianza en cada extracción.

\subsubsection{Mitigación de alucinaciones}

Dado que la alucinación en \acs{LLM} es una limitación innata que 
solo puede mitigarse y no eliminarse~\cite{xu2024hallucination}, 
el sistema adopta un esquema de verificación en capas que 
reduce su incidencia y preserva la trazabilidad necesaria para 
que el analista identifique los errores que escapen al sistema.

La verificación automática se aplica antes de pasar a la fase 
de agregación. La estrategia se adapta a la naturaleza del 
dato, distinguiendo entre valores literales y valores 
inferidos, ya que cada uno admite un tipo distinto de 
contraste contra la fuente.


Los valores literales, como fechas, nombres y montos, se 
verifican mediante coincidencia de texto contra el fragmento 
indicado por su fuente. Esta verificación detecta directamente 
las fabricaciones del modelo a un costo computacional mínimo, 
en línea con técnicas de auto-verificación 
factual~\cite{manakul2023selfcheckgpt}. Si el valor aparece 
en el fragmento, se considera verificado. En caso contrario, 
se marca como sospechoso y se envía a una nueva consulta al 
modelo, que recibe el dato junto al fragmento y debe 
reextraerlo o confirmar su ausencia. El resultado reemplaza 
al dato inicial.

Los valores inferidos, que el modelo deduce a partir del 
texto en lugar de tomarlos literalmente, no admiten una 
verificación por coincidencia. Para estos casos adoptamos un 
esquema de verificación encadenada 
(\textit{chain-of-verification})~\cite{dhuliawala2024chainofverification}, 
en el que el modelo revisa su propia respuesta contra el 
fragmento original y la justificación que produjo, y confirma 
si la inferencia se desprende razonablemente del texto. Si 
falla, el dato se reextrae añadiendo el contexto explícito.

La combinación de ambos mecanismos responde a un compromiso 
entre costo computacional y capacidad de razonamiento. La 
comparación de texto resuelve los casos simples sin invocar 
al modelo, lo que reduce el costo y elimina la posibilidad 
de alucinación en datos verificables sintácticamente. La 
verificación encadenada se reserva para los datos que 
requieren interpretación, donde el costo adicional se 
justifica por el riesgo de error. Esta selectividad reduce 
la tasa global de alucinaciones sin penalizar el tiempo de 
procesamiento.

Aun así, el resultado de la verificación automática no 
constituye una garantía. Por la limitación teórica antes 
señalada, una fracción de errores puede atravesar ambos 
filtros, especialmente en casos donde la evidencia textual es 
ambigua o donde el modelo refuerza su propia interpretación 
errónea durante la verificación encadenada. Por ello, los 
tres anclajes producidos durante la extracción semántica, 
fuente, fragmento de evidencia y justificación, se conservan 
en la representación final del informe. La interfaz expone 
estos anclajes en la fase de exploración, de modo que el 
analista pueda contrastar cada dato con su contexto original 
y emitir un juicio informado. Esta separación de 
responsabilidades, donde el sistema reduce la incidencia de 
errores y el analista valida el resultado, es coherente con 
los marcos de análisis visual que conciben la interpretación 
como una tarea conjunta entre el cómputo automático y el 
razonamiento humano~\cite{keim2008visual, sacha2014}.

\subsection{Agregación}

La fase de agregación toma los informes procesados 
individualmente en la etapa anterior y los consolida en una 
representación unificada por caso. El reto central de esta 
fase es la resolución de correferencia entre informes, ya que 
una misma variable puede aparecer descrita con texto diferente 
en distintos informes del mismo caso, y sin un mecanismo de 
resolución los datos no pueden integrarse en una secuencia 
coherente~\cite{cattan2021crossdoc}. La fase se compone de 
cuatro pasos que se ejecutan en orden, agrupación por caso, 
emparejamiento semántico de variables, verificación cruzada y 
construcción del vector de características.

\subsubsection{Agrupación por caso}

Los informes que comparten un mismo caso se agrupan, 
conformando su secuencia ordenada en el tiempo. Los datos 
generales del caso, como el título o la descripción, se fusionan 
en una representación única tomando como base el primer informe 
disponible y completando los campos ausentes con los informes 
posteriores. Esta política de fusión refleja el supuesto de que 
los informes posteriores complementan, antes que contradicen, 
la información del primero. Las contradicciones reales se 
detectan en el paso de verificación cruzada.

\subsubsection{Emparejamiento semántico de variables}

Una vez agrupados los informes de un caso, el sistema compara 
las variables entre informes consecutivos para identificar 
cuáles representan el mismo elemento descrito con texto 
diferente. Este emparejamiento se realiza mediante una consulta 
al \acs{LLM}, que recibe las variables de informes consecutivos junto 
con su justificación y determina cuáles corresponden al mismo 
elemento evolucionado, registrando el cambio observado entre 
informes.

La decisión de delegar el emparejamiento al modelo en lugar de 
basarlo únicamente en similitud léxica responde a la naturaleza 
del texto institucional. Una misma variable puede reformularse 
de un informe al siguiente, cambiar de tono o incorporar nuevos 
detalles, sin que ello implique un cambio en el referente. 
Métricas léxicas tradicionales, como similitud de cadenas o 
\textit{embeddings} de palabra, no capturan adecuadamente estos casos 
cuando la reformulación es sustantiva. Trabajos recientes en 
consistencia entre resúmenes han mostrado que delegar la 
comparación a un modelo de lenguaje produce decisiones más 
robustas frente a paráfrasis y reformulaciones~\cite{laban2022summac}, 
y es ese el tipo de robustez que el dominio requiere.

Durante este proceso, las variables que no aparecen en todos 
los informes se registran con valor nulo, distinguiendo 
explícitamente entre un dato ausente y un dato con valor cero. 
Esta distinción es necesaria porque, en un caso, la ausencia de 
una variable en un informe específico es una señal informativa 
en sí misma, no un error de extracción.

\subsubsection{Verificación cruzada}

Con el caso ya reconstruido, el sistema revisa inconsistencias 
entre informes que no eran visibles al procesar cada documento 
por separado. Esto incluye fechas contradictorias, nombres 
distintos para un mismo actor y hallazgos cuyo estado declarado 
no es consistente con su evolución a lo largo del 
caso~\cite{manakul2023selfcheckgpt}.

Este paso es complementario al de mitigación de alucinaciones 
de la fase anterior. Mientras aquella verificación contrasta 
cada dato contra su fragmento fuente, la verificación cruzada 
contrasta cada dato contra el resto de los informes del caso. 
Capturan tipos de error distintos. Una fecha extraída 
correctamente de un fragmento puede entrar en conflicto con la 
fecha extraída correctamente de otro fragmento si los documentos 
mismos no coinciden, y solo una verificación a nivel de caso 
detecta ese tipo de inconsistencia.
%¿por qué no usaron el LLM para resolver?", tu respuesta tiene que ser exactamente la del segundo párrafo. Lo intentamos así inicialmente, pero el modelo terminaba inventando razones para preferir uno u otro valor sin que existiera evidencia textual real para hacerlo. Eso es contradictorio con el objetivo de mitigar alucinaciones. La decisión de marcar y exponer es más conservadora, más honesta, y además aprovecha el valor informativo del conflicto.
Cuando se detecta una inconsistencia, el sistema no la resuelve 
automáticamente. Conserva los valores en conflicto junto con 
sus respectivos anclajes al documento fuente y los marca como 
inconsistentes en la representación final del caso. Esta 
decisión responde a dos consideraciones. Por un lado, una 
resolución automática mediante una nueva consulta al modelo 
introduce el mismo riesgo de alucinación que el sistema busca 
mitigar, ya que el modelo podría producir una respuesta 
plausible sin que existan elementos para preferir un valor 
sobre el otro. Por otro, una inconsistencia entre informes de 
un mismo caso es a menudo informativa en sí misma. Cuando dos 
documentos del mismo expediente reportan estados contradictorios 
sobre un mismo hallazgo, esa contradicción suele ser parte del 
fenómeno bajo análisis, no un error de extracción que conviene 
resolver. Marcar la inconsistencia y exponerla al analista 
preserva esa señal y respeta la división de roles entre el 
sistema y el analista propia del análisis 
visual~\cite{keim2008visual}.

Solo se aplica una normalización determinística para 
inconsistencias puramente sintácticas, como variantes de 
formato de fecha o diferencias de mayúsculas y acentos en 
nombres propios, donde la equivalencia es decidible sin 
interpretación.

\subsubsection{Construcción del vector de características}

Por último, el sistema calcula un vector de características 
para cada caso. Este vector combina dos componentes 
heterogéneos. El primero es un conjunto de métricas numéricas 
derivadas de las variables del caso, como la proporción de 
hallazgos de Seguimiento cerrados, el porcentaje de avance de 
las variables de Progreso, la proporción de condiciones de 
Verificación cumplidas y el número de actores involucrados. El 
segundo es un \textit{embedding} semántico generado a partir 
de las conclusiones registradas en los informes del caso. Para 
el \textit{embedding} se utiliza un modelo de \textit{Sentence 
Transformers}~\cite{reimers2019sbert}, que proyecta el texto a 
un espacio vectorial denso preservando su significado semántico.

La decisión de combinar componentes numéricos y semánticos en 
un solo vector responde al objetivo de soportar las dos 
formas en que un analista puede juzgar la similitud entre 
casos. Dos casos pueden parecerse porque sus indicadores 
cuantitativos coinciden, o porque sus conclusiones narrativas 
describen situaciones análogas. Un vector que ignore una de 
las dos dimensiones colapsaría una distinción que el analista 
sí percibe. La combinación obliga, no obstante, a normalizar 
ambos componentes antes de concatenarlos para que ninguno 
domine la métrica de distancia, ajuste que se realiza mediante 
escalado por desviación estándar.

El vector resultante es el insumo de la vista de similitud, 
que lo proyecta a dos dimensiones mediante 
\ac{UMAP}~\cite{mcinnes2018umap}. Optamos por \acs{UMAP} frente a 
alternativas como t-SNE porque preserva mejor la estructura 
global del conjunto de datos, lo que es relevante para una 
tarea de exploración donde la posición relativa entre grupos 
de casos comunica información al analista, no solo la 
vecindad local entre puntos individuales.

Al finalizar esta fase, el sistema produce tres archivos por 
conjunto de casos analizado. Un consolidado por caso con todos 
sus informes y variables sincronizadas. Un archivo de 
coordenadas geográficas por caso. Y un archivo de vectores de 
características, uno por caso. Estos archivos constituyen la 
entrada de la fase de exploración interactiva.

\subsection{Exploración interactiva}

La fase de exploración interactiva expone la representación 
unificada producida por las fases anteriores al analista. Su 
diseño parte del principio de vistas múltiples coordinadas, 
que sostiene que un fenómeno multidimensional se comprende 
mejor cuando cada dimensión se representa con la 
visualización más adecuada y las vistas se vinculan mediante 
interacción 
sincronizada~\cite{wang2000guidelines, roberts2007cmv}.

\begin{figure}[htbp]
    \centering
    \includegraphics[width=\textwidth]{vista_p.png}
    \caption{Interfaz general del sistema sobre un caso de 
    estudio del Proceso Electoral 2026, donde cuatro Oficinas 
    Descentralizadas de Procesos Electorales se analizan en 
    paralelo. La interfaz se compone de cuatro regiones 
    coordinadas, la vista geoespacial (\textbf{A}), la vista 
    temporal (\textbf{B}), la vista de similitud (\textbf{C}) 
    y la vista de detalle (\textbf{D}).}
    \label{fig:vista_principal}
\end{figure}

La exploración sigue un flujo de cuatro etapas. El analista 
parte de una vista panorámica del conjunto de casos, ya sea 
sobre el territorio en la vista geoespacial (\textbf{A}) o 
sobre el espacio semántico de similitud (\textbf{B}). Al 
seleccionar uno o varios casos en cualquiera de las dos, la 
selección se propaga a la otra y se construye la vista 
temporal (\textbf{C}) sobre los casos elegidos. Desde la 
vista temporal, el analista puede inspeccionar un informe, 
una variable o una región de la línea de tiempo, lo que 
abre la vista de detalle (\textbf{D}) con la información 
asociada. Y desde la vista de detalle, los enlaces a los 
fragmentos fuente abren la vista de evidencia, que muestra 
el documento original con el fragmento resaltado. Un panel 
de filtros adaptativos en el encabezado complementa este 
flujo, permitiendo restringir la exploración por tiempo, 
espacio y características del caso.

A continuación, detallamos cada componente siguiendo el recorrido del usuario: desde las primeras vistas de exploración hasta el acceso a la evidencia original.

\subsubsection{Vista geoespacial}

La vista geoespacial (\textbf{A} en la 
Figura~\ref{fig:vista_principal})~localiza los casos sobre 
un mapa interactivo y soporta la tarea T3. Aprovecha que 
cada caso tiene una ubicación asociada, calculada en la 
fase de agregación a partir de las referencias geográficas 
extraídas de los informes.

La vista geoespacial soporta selección de 
casos individuales o de regiones territoriales. La 
selección se propaga a la vista de similitud, donde los 
puntos correspondientes se resaltan, y a la vista 
temporal, que se construye sobre el subconjunto elegido.


\subsubsection{Vista de similitud}

La vista de similitud (\textbf{B} en la 
Figura~\ref{fig:vista_principal})~proyecta el conjunto 
completo de casos en un espacio bidimensional para 
soportar la tarea T4, donde el analista busca identificar 
grupos de casos parecidos o casos atípicos respecto al 
conjunto. Su fuente es el vector de características 
construido en la fase de agregación, que combina 
indicadores numéricos con un \textit{embedding} semántico 
de las conclusiones del caso.

La proyección se realiza con \acs{UMAP}~\cite{mcinnes2018umap}, 
cuya elección frente a alternativas como t-SNE ya fue 
discutida en la fase de agregación. La representación es 
un gráfico de dispersión donde cada punto corresponde a 
un caso. La cercanía espacial entre puntos indica que los 
casos comparten un perfil cuantitativo y narrativo 
similar, lo cual permite al analista detectar agrupaciones 
por su sola geometría sin tener que inspeccionarlos uno 
por uno.

La vista de similitud es una entrada complementaria a la 
geoespacial. Mientras la geoespacial agrupa casos por 
proximidad territorial, la de similitud los agrupa por 
proximidad semántica y cuantitativa. Esta dualidad permite 
que el analista descubra, por ejemplo, que dos casos 
geográficamente distantes presentan trayectorias casi 
idénticas, o que dos casos vecinos difieren en su 
comportamiento. Al posicionar el cursor sobre un punto, 
una ventana flotante muestra el resumen del caso. La 
selección de un punto, o de un grupo de puntos mediante 
\textit{lasso}, propaga el contexto a la vista 
geoespacial y construye la vista temporal sobre los casos 
elegidos.

\subsubsection{Filtros adaptativos}

El encabezado de la interfaz expone un panel de filtros 
que opera sobre la selección activa. Tres dimensiones de 
filtrado están disponibles. La dimensión temporal, 
mediante un control deslizante que limita la ventana de 
fechas considerada. La dimensión espacial, mediante 
selectores territoriales que restringen la región 
analizada. Y la dimensión de características, mediante 
controles que filtran por atributos de las variables, 
como el estado de una variable de Seguimiento o el rango 
de valores de una variable de Progreso.

Los controles de la dimensión de características se 
adaptan al tipo de variable activa. Esta adaptación responde a un 
principio elemental de diseño de interfaz, según el cual 
los controles deben corresponder a la estructura de los 
datos sobre los que 
operan~\cite{shneiderman1996tasks}, y evita el patrón 
común de exponer controles que no aplican a la variable 
activa, que confunde al usuario y degrada la experiencia 
de exploración.

\subsubsection{Vista temporal}

La vista temporal (\textbf{C} en la 
Figura~\ref{fig:vista_principal})~se construye sobre los 
casos seleccionados en las vistas de entrada y soporta 
las tareas T1 y T2. Su representación se adapta al tipo 
de variable bajo análisis, siguiendo el principio de que 
una sola codificación visual no puede expresar 
adecuadamente la heterogeneidad de los seis tipos del 
modelo de datos.

Por defecto, cuando ninguna variable específica está 
seleccionada, la vista temporal muestra los informes de 
los casos como un diagrama de Gantt, donde el eje 
horizontal corresponde al tiempo y el eje vertical 
organiza los casos en filas paralelas. Esta 
representación responde a una propiedad estructural del 
dominio. Los informes institucionales se caracterizan 
por intervalos con duración explícita, como un plazo de 
respuesta o una ventana de ejecución, intercalados con 
eventos puntuales, como la emisión del informe mismo. 
La sintaxis del Gantt preserva ambas clases y permite 
identificar visualmente retrasos, periodos sin actividad 
y solapamientos entre casos. La organización en filas 
paralelas, además, soporta la comparación directa entre 
casos sin requerir una segunda vista.

Al seleccionar una variable de \textbf{Seguimiento}, la 
vista temporal cambia a un diagrama de Sankey, donde 
cada flujo representa la transición de un hallazgo entre 
estados a lo largo de los informes. Esta elección 
responde a la estructura semántica del tipo. Las 
variables de Seguimiento describen elementos que se 
resuelven, persisten o generan nuevas consecuencias 
entre informes, lo cual es exactamente el patrón que 
los diagramas de Sankey representan, flujos entre 
estados con grosor proporcional al volumen del 
flujo~\cite{riehmann2005sankey}. La representación 
permite que el analista identifique, por ejemplo, qué 
proporción de hallazgos abiertos en un informe inicial 
quedó cerrada en un informe posterior, y qué proporción 
persistió o derivó en nuevos hallazgos.

Al seleccionar una variable de \textbf{Progreso}, la 
vista temporal cambia a un diagrama de líneas, donde una 
línea representa el valor observado a lo largo del 
tiempo y una segunda línea, en estilo discontinuo, 
representa el valor estimado o meta. Esta es la 
codificación canónica para series temporales numéricas 
con referencia~\cite{cleveland1985elements} y permite 
detectar visualmente desvíos respecto a la meta sin 
recurrir a cálculos auxiliares.

Para los cuatro tipos restantes, la vista temporal adopta 
codificaciones diferenciadas según la estructura semántica 
de cada uno. Una variable de \textbf{Estado}, que asigna 
una categoría a cada informe, se representa como una 
banda horizontal segmentada por color, siguiendo el 
patrón de las visualizaciones de tipo \textit{swimlane} 
que codifican secuencias categóricas a lo 
largo del tiempo~\cite{wongsuphasawat2011lifeflow}. Una 
variable \textbf{Relacional}, que describe los actores 
involucrados y sus roles, se representa como un grafo 
dinámico cuando el conjunto de actores es pequeño, o 
como una matriz de adyacencia ordenada cronológicamente 
cuando es 
denso~\cite{beck2017taxonomy}. Una variable de tipo 
\textbf{Verificación}, que registra el cumplimiento 
booleano de una condición a lo largo de los informes, se 
representa como un \textit{dot plot} binario donde cada 
fila corresponde a un caso y cada columna a un informe, 
codificación que permite identificar rápidamente patrones 
de incumplimiento sistemático. Las variables de tipo 
\textbf{Narrativo}, finalmente, no admiten una 
codificación temporal compacta porque su contenido es 
texto extenso. En su lugar, se exponen como tarjetas en 
la vista de detalle al seleccionar el informe 
correspondiente, manteniendo el acceso al texto sin 
forzar una representación visual que perdería 
información.

\subsubsection{Vista de detalle}

La vista de detalle (\textbf{D} en la 
Figura~\ref{fig:vista_principal})~se activa al 
seleccionar un informe, un caso o una región temporal en 
las otras vistas. Presenta la información asociada a esa 
selección, organizada en bloques que cubren los 
metadatos del elemento, las variables registradas en él 
y los enlaces a sus fragmentos fuente.

Para un informe seleccionado, la vista muestra el título 
del hito, las fechas relevantes (emisión del informe, 
plazo de respuesta), las situaciones identificadas y la 
referencia al documento original. Cada dato visible en 
esta vista es interactivo y abre la vista de evidencia 
sobre el fragmento correspondiente del documento.

Esta vista funciona como puente entre las 
representaciones agregadas y la evidencia textual. Las 
vistas A, B y C operan sobre representaciones derivadas 
del proceso de extracción, y aunque son útiles para 
reconocer patrones, no exponen el texto original que 
sustenta cada dato. La vista de detalle traduce un 
elemento seleccionado en sus componentes verificables, y 
expone los enlaces que permiten descender al documento 
fuente.

\subsubsection{Vista de evidencia}

La vista de evidencia materializa los anclajes producidos 
durante la extracción y soporta la tarea T5. Cuando el 
analista activa cualquier enlace desde la vista de 
detalle, la vista de evidencia se abre como una ventana 
embebida sobre la interfaz, carga el documento original, 
navega a la página correspondiente y resalta el 
fragmento exacto del que se obtuvo el dato. La 
Figura~\ref{fig:vista_pdf} muestra esta interacción 
sobre un informe de auditoría, donde el recuadro 
amarillo señala el fragmento fuente del dato 
seleccionado.

\begin{figure}[htbp]
    \centering
    \includegraphics[width=\textwidth]{vista_pdf.png}
    \caption{Vista de evidencia. El sistema abre el 
    documento original, navega a la página correspondiente 
    y resalta el fragmento desde el cual se extrajo el 
    dato seleccionado.}
    \label{fig:vista_pdf}
\end{figure}

La vista ofrece tres elementos de soporte para la 
verificación. Primero, la navegación entre páginas del 
documento permite consultar el contexto previo y 
posterior al fragmento, ya que un fragmento aislado 
puede no ser suficiente para juzgar la corrección de la 
extracción. Segundo, el resaltado visual del fragmento 
fuente reduce el costo cognitivo de localizarlo dentro 
de la página, especialmente en informes densos donde el 
fragmento relevante puede ser una sola línea entre 
cientos. Tercero, junto al fragmento se expone la 
justificación que el modelo articuló al realizar la 
extracción, lo que permite al analista contrastar no 
solo el dato sino también el razonamiento que lo 
produjo. Esta vista es central al diseño del sistema. Las fases 
anteriores adoptaron una posición explícita sobre la 
limitación teórica de los \acs{LLM} frente a la 
alucinación, asumiendo que el sistema mitiga pero no 
elimina los errores. Esa posición sería declarativa sin 
un mecanismo concreto que permita al analista verificar 
cada dato. La vista de evidencia es ese mecanismo. Al 
colocar el documento original a un clic de cualquier 
dato extraído, el sistema se alinea con la noción de calibración 
de confianza reportada en la literatura~\cite{yeh2025storyribbons}.
